% TOP_PROBLEM: {"id": 2305055, "title": "Research Problems in Function Theory — Problem 5.55", "category_id": 8, "category": {"id": 8, "name": "analysis", "display_name": "Analysis", "description": "Limits, continuity, calculus, and function theory.", "slug": "analysis", "order_index": 8, "created_at": "2026-07-31T15:26:25.671Z"}, "status": "open"}
% FIRST_POSTED: null
% ENTRY_KIND: statement_only
% Imported from: batch01.tex; UnsolvedMath ID 2305055
\documentclass[11pt]{book}
\usepackage{fontspec}
\setmainfont{Latin Modern Roman}
\setsansfont{Latin Modern Sans}
\usepackage{amsmath,amssymb,mathtools}
\usepackage{unicode-math}
\setmathfont{Latin Modern Math}
\usepackage{xurl}
\usepackage[hidelinks]{hyperref}
\newcommand{\sref}[2]{\hyperlink{source:#1:#2}{[S#2]}}
\newcommand{\cataloguescope}[1]{\paragraph{Mathematical question.}#1}
\begin{document}
% UnsolvedMath ID 2305055; source catalogue code AMR-022-5055
\setcounter{section}{2305054}
\section{Research Problems in Function Theory — Problem 5.55}\label{2305055}
{\small\sffamily UnsolvedMath ID 2305055 · Analysis}
\subsection{Definitions and mathematical statement}

\paragraph{Shared notation.}$\mathbb N=\{1,2,\ldots\}$, and $\mathbb Z,\mathbb Q,\mathbb R,\mathbb C$ denote the integers, rationals, reals and complex numbers. Any other conventions are specified locally.


Let $\mathbb D=\{z\in\mathbb C:|z|<1\}$ and $\mathbb T=\partial\mathbb D$. For an analytic $h$ put $G_h=h(\mathbb D)$. Say $h$ has the $L$-property on a sequence $(z_j)$ with $|z_j|\to1$ if, for every compact $K\subset G_h$, only finitely many indices $j$ satisfy $h(z_j)\in K$. An ordered $L$-pair $(f,g)$ means that every such sequence on which $f$ has the $L$-property is also one on which $g$ has it. A nonconstant analytic $\alpha$ is an $L$-atom if, for every analytic $f$ for which $(f,\alpha)$ is an ordered $L$-pair, there is an analytic $\phi_f:G_\alpha\to\mathbb C$ with $f=\phi_f\circ\alpha$. For a general planar domain $D$, replace $|z_j|\to1$ by the requirement that $(z_j)$ eventually leaves every compact subset of $D$, and use the same range-compactness and factorization definitions.
\paragraph{Statement recorded in the supplied source.}
Is $\alpha$ an $L$-atom if and only if $\alpha$ is univalent (injective)? Determine what becomes of this equivalence when the source domain is a general planar domain $D$. \sref{2305055}{2}
\subsection{Short English statement}
Characterize the analytic functions that force every preceding member of an ordered escape pair to factor through them; are these exactly the univalent functions, in the disk and in other domains?
\subsection{Sources}
\begin{description}
\item[\hypertarget{source:2305055:1}{S1}] ULAM.AI and the UnsolvedMath Contributors, \emph{UnsolvedMath: A Curated Collection of Open Mathematics Problems}, record 2305055 (AMR-022-5055). CC BY 4.0. \url{https://huggingface.co/datasets/ulamai/UnsolvedMath}.
\item[\hypertarget{source:2305055:2}{S2}] Walter K. Hayman and Edward F. Lingham, \emph{Research Problems in Function Theory} (2018), arXiv:1809.07200. Problem 5.55, complete original question, all following updates, and the relevant chapter definitions. \url{https://arxiv.org/abs/1809.07200}.
\end{description}
\label{end:2305055}
\end{document}
